
\documentclass{article}
\usepackage{colortbl}
\usepackage{makecell}
\usepackage{multirow}
\usepackage{supertabular}

\begin{document}

\newcounter{utterance}

\centering \large Interaction Transcript for game `wordle\_withcritic\_noexp', experiment `medium\_frequency\_words\_clue\_with\_critic', episode 17 with reflexion{-}Qwen3.5{-}9B.
\vspace{24pt}

{ \footnotesize  \setcounter{utterance}{1}
\setlength{\tabcolsep}{0pt}
\begin{supertabular}{c@{$\;$}|p{.15\linewidth}@{}p{.15\linewidth}p{.15\linewidth}p{.15\linewidth}p{.15\linewidth}p{.15\linewidth}}
    \# & $\;$A & \multicolumn{4}{c}{Game Master} & $\;\:$B\\
    \hline

    \theutterance \stepcounter{utterance}  
    & & \multicolumn{4}{p{0.6\linewidth}}{
        \cellcolor[rgb]{0.9,0.9,0.9}{
            \makecell[{{p{\linewidth}}}]{
                \texttt{\tiny{[P1$\langle$GM]}}
                \texttt{You have attempted this same task before and did not succeed. These are the plans you wrote after those attempts.} \\
\texttt{Trial \#1: New plan: I failed because I generated a response containing only the "guess:" tag without the required "explanation:" tag, violating the format rules. Next time, I will strictly adhere to the required format by providing a concise explanation followed by the guess on separate lines, ensuring no extra text or missing tags are included. For the first turn, I will explain my strategy of using the clue "virile" to select a word with strong consonants and common vowels, then provide the guess.} \\
\\ 
\texttt{Trial \#2: New plan: I failed because I included an "explanation:" tag in my response, but the Game Master explicitly instructed me to "Do not generate any other text" and only use the "guess:" tag. The rules did not require an explanation, only the guess itself. Next time, I will strictly output only the "guess:" tag followed by the word, omitting any explanatory text or additional tags to ensure the response is well{-}formed and legal. For the first turn, I will directly provide a guess related to the clue "virile" without any preamble.} \\
\\ 
\texttt{Trial \#3: New plan: I failed because I violated the strict formatting rule "Do not generate any other text" by including an "explanation:" tag, which the Game Master explicitly forbade in this specific instance (unlike Example 2 where explanations were required). The rules for this game only mandate the "guess:" tag. Next time, I will strictly output only the "guess:" tag followed by a valid five{-}letter word that matches the clue "virile" (e.g., "manly" or "strong"), ensuring no explanatory text, no extra tags, and no preamble are included. For the first turn, I will immediately provide: guess: manly.} \\
\\ 
\\ 
\texttt{You are a language wizard who likes to guess words by using the given rules.} \\
\\ 
\texttt{Welcome to Wordle! You have six attempts to guess the target word, a valid English word of five lowercase letters (a{-}z). Please use the tag "guess:" to provide your guess.} \\
\\ 
\texttt{To help you make an informed guess, you will receive a clue for the word, such as} \\
\texttt{clue: "snowy white"} \\
\\ 
\texttt{Here is an example guess based on the clue: } \\
\texttt{guess: apple} \\
\\ 
\texttt{I will then indicate whether I agree or disagree with your guess, but agreeing with a guess does not confirm its correctness. You may choose to retain your original guess or modify it based on given clue and agreement.} \\
\\ 
\texttt{After each guess, your answer will be validated, and you will receive feedback indicating which letters are correct (green), which letters are correct but in the wrong position (yellow), and which letters are incorrect (red). This feedback can be useful in determining which letters to include or exclude in your next guess.} \\
\\ 
\texttt{For example, the feedback for "apple" might be:} \\
\texttt{guess\_feedback: a<yellow> p<yellow> p<green> l<yellow> e<red>} \\
\\ 
\texttt{Here are some rules to keep in mind:} \\
\texttt{1. Please do not generate any guess feedback or clue on your own} \\
\texttt{2. Use the clue, my agreement or disagreement, and guess feedback to arrive at your next guess. } \\
\\ 
\texttt{Do not generate any other text.} \\
\\ 
\texttt{Let's begin with your first guess.} \\
\\ 
\texttt{clue: virile} \\
            }
        }
    }
    & \\ \\

    \theutterance \stepcounter{utterance}  
    & \multicolumn{4}{p{0.6\linewidth}}{
        \cellcolor[rgb]{0.9,0.9,0.9}{
            \makecell[{{p{\linewidth}}}]{
                \texttt{\tiny{[P1$\rangle$GM]}}
                \texttt{guess: manly} \\
            }
        }
    }
    & & \\ \\

    \theutterance \stepcounter{utterance}  
    & & \multicolumn{4}{p{0.6\linewidth}}{
        \cellcolor[rgb]{0.9,0.9,0.9}{
            \makecell[{{p{\linewidth}}}]{
                \texttt{\tiny{[GM$\rangle$P2]}}
                \texttt{I need your assistance with a word game in which we need to find a 5{-}letter word using a clue and a guess.} \\
\\ 
\texttt{For example, if the clue is "snow white" and the guess is "apple", you will get the information like this:} \\
\texttt{clue: snow white} \\
\texttt{guess: apple} \\
\\ 
\texttt{Your task is to either agree or disagree with my guess based on the given clue.} \\
\\ 
\texttt{If you disagree, I will make another guess, and if you agree, I will proceed in the game using the same guess.} \\
\\ 
\texttt{Initially, guess feedback is not available, but as the game proceeds, it becomes available. The guess feedback is not for current guess, but rather an indication of what letters can be present in the current guess. A letter highlighted in green indicates that it is a correct letter in the correct position, while yellow indicates that it is a correct letter in the wrong position, and red indicates a wrong letter. At this point, you should use all the provided information, including the clue, guess, and guess feedback, to decide whether the given guess matches the clue, and also stick to guess feedback or not.} \\
\\ 
\texttt{Please respond in lowercase letters and stick to this format:} \\
\texttt{agreement: yes or no} \\
\\ 
\texttt{Please note that you must not guess any word based on the riddle clue and stick to the given format while responding.} \\
\\ 
\texttt{Do not generate any other text.} \\
\\ 
\texttt{Let's begin!} \\
\\ 
\texttt{clue: virile} \\
\texttt{guess: manly} \\
            }
        }
    }
    & \\ \\

    \theutterance \stepcounter{utterance}  
    & & & \multicolumn{4}{p{0.6\linewidth}}{
        \cellcolor[rgb]{0.9,0.9,0.9}{
            \makecell[{{p{\linewidth}}}]{
                \texttt{\tiny{[GM$\langle$P2]}}
                \texttt{agreement: yes} \\
            }
        }
    }
     \\ \\

    \theutterance \stepcounter{utterance}  
    & & \multicolumn{4}{p{0.6\linewidth}}{
        \cellcolor[rgb]{0.9,0.9,0.9}{
            \makecell[{{p{\linewidth}}}]{
                \texttt{\tiny{[P1$\langle$GM]}}
                \texttt{clue: virile} \\
\texttt{guess\_agreement: yes} \\
\texttt{Provide your response only in this format.} \\
\texttt{guess: word} \\
            }
        }
    }
    & \\ \\

    \theutterance \stepcounter{utterance}  
    & \multicolumn{4}{p{0.6\linewidth}}{
        \cellcolor[rgb]{0.9,0.9,0.9}{
            \makecell[{{p{\linewidth}}}]{
                \texttt{\tiny{[P1$\rangle$GM]}}
                \texttt{guess: manly} \\
            }
        }
    }
    & & \\ \\

    \theutterance \stepcounter{utterance}  
    & & & \multicolumn{2}{p{0.3\linewidth}}{
        \cellcolor[rgb]{0.9,0.9,0.9}{
            \makecell[{{p{\linewidth}}}]{
                \texttt{\tiny{[GM$|$GM]}}
                \texttt{attempts = 1} \\
\texttt{target\_word = manly} \\
\texttt{target\_word\_clue = virile} \\
\texttt{guess = manly} \\
\texttt{guess\_feedback = m<green> a<green> n<green> l<green> y<green>} \\
            }
        }
    }
    & & \\ \\

    \theutterance \stepcounter{utterance}  
    & & & \multicolumn{2}{p{0.3\linewidth}}{
        \cellcolor[rgb]{0.9,0.9,0.9}{
            \makecell[{{p{\linewidth}}}]{
                \texttt{\tiny{[GM$|$GM]}}
                \texttt{game\_result = WIN} \\
            }
        }
    }
    & & \\ \\

\end{supertabular}
}

\end{document}
